\chapter{Abscesses}
\index{Abscesses}

\section*{Definition and Overview}
An abscess is a painful, localised collection of pus that forms when the immune system walls off an infection. Abscesses can develop almost anywhere in the body: beneath the skin, in the mouth, around the anus, or within internal organs such as the liver, lungs, and brain. Because the pus is trapped inside a fibrous cavity, an abscess rarely resolves with antibiotics alone and usually requires drainage. This chapter covers the causes, presentation, diagnosis, and management of abscesses as a clinical problem; the technique of draining them is described in the chapter on abscess incision and drainage.

\section*{Epidemiology}
Abscesses are common, and skin and soft-tissue abscesses are among the most frequent problems seen in emergency departments. Deep abscesses, affecting the abdomen, liver, lung, or brain, are less common and typically occur in people with specific risk factors. Higher-risk groups include people with diabetes, weakened immune systems, or poor local blood supply, people who inject drugs, and people recovering from surgery or penetrating injuries. Abscesses affect all ages; some types, such as perianal abscesses, are more common in men.

\section*{Aetiology and Pathophysiology}
Abscesses are caused by bacteria, and occasionally by parasites or fungi, that provoke a vigorous inflammatory response. White blood cells migrate to the site of infection, and the accumulating mixture of dead cells, bacteria, and tissue debris forms pus. The body then surrounds the pus with a tough wall of fibrous tissue. This wall protects surrounding structures but also seals the infection away from antibiotics and immune cells, so the pressure within the cavity rises and the abscess cannot clear of its own accord.

\begin{itemize}
    \item \textbf{Skin:} \textit{Staphylococcus aureus}, including MRSA, and streptococci
    \item \textbf{Mouth and teeth:} mixed oral flora, including anaerobes
    \item \textbf{Bowel and anus:} intestinal bacteria, including anaerobes and enteric organisms
    \item \textbf{Liver:} organisms ascending the bile duct or carried in the blood, and occasionally amoeba
    \item \textbf{Lung:} bacteria aspirated from the mouth and throat
    \item \textbf{Brain:} organisms carried in the blood or spreading from the sinuses, ears, or teeth
\end{itemize}

\section*{Clinical Features}
The clinical picture depends on the location of the abscess.

\begin{itemize}
    \item \textbf{Skin:} a hot, red, swollen, painful lump that becomes soft and fluctuant over days
    \item \textbf{General:} fever, chills, malaise, and a raised heart rate when infection is spreading
    \item \textbf{Dental:} severe toothache, facial swelling, and pain on chewing
    \item \textbf{Throat (peritonsillar):} severe one-sided sore throat, difficulty swallowing, and a muffled or "hot potato" voice
    \item \textbf{Liver:} right upper abdominal pain, fever, nausea, and sometimes jaundice
    \item \textbf{Lung:} cough with foul-tasting sputum, fever, and pleuritic chest pain
    \item \textbf{Brain:} headache, fever, and neurological symptoms such as weakness, seizures, or confusion
\end{itemize}

\section*{Differential Diagnosis}
Several conditions can resemble an abscess and should be distinguished from it.

\begin{itemize}
    \item A solid tumour, which may mimic a superficial or deep abscess
    \item A haematoma or seroma, collections of blood or fluid following injury or surgery
    \item An enlarged lymph node, including from tuberculosis or lymphoma
    \item Cellulitis without a drainable collection
    \item A cyst or benign lump, such as an epidermoid cyst
    \item Deep vein thrombosis or inflammatory bowel disease, depending on the site
\end{itemize}

\section*{When to Seek Medical Care}
Any lump that may be an abscess should be assessed by a doctor, especially if it is red, hot, painful, or enlarging.

\textbf{Contact your local emergency services for:}
\begin{itemize}
    \item Rapidly spreading redness or swelling
    \item High fever with chills, confusion, or drowsiness, which may indicate sepsis
    \item Difficulty breathing or swallowing with a throat, mouth, or neck abscess
    \item Severe headache, neck stiffness, or neurological symptoms with fever, which may indicate a brain abscess
\end{itemize}

\section*{Diagnosis and Investigations}
Most superficial abscesses are diagnosed clinically, while deep abscesses require imaging.

\begin{itemize}
    \item \textbf{Clinical examination:} history and physical findings are often sufficient for superficial collections
    \item \textbf{Ultrasound:} confirms superficial abscesses and some deep collections
    \item \textbf{CT or MRI:} locates and characterises abscesses in the abdomen, chest, or brain
    \item \textbf{Blood tests:} full blood count and inflammatory markers, with blood cultures in systemically unwell patients
    \item \textbf{Microbiological culture:} pus obtained by aspiration or drainage is sent to identify the causative organism
    \item \textbf{Imaging-guided aspiration:} a needle sample may be obtained for both diagnosis and treatment
\end{itemize}

\section*{Management}
Drainage is the cornerstone of treatment, with antibiotics and treatment of the source as supporting measures.

\begin{itemize}
    \item \textbf{Drainage:} incision and drainage for skin abscesses; image-guided catheter drainage or surgery for deep abscesses
    \item \textbf{Antibiotics:} usually combined with drainage for cellulitis, deep abscesses, or systemic illness, and for people with diabetes or weakened immunity
    \item \textbf{Treat the source:} such as extracting an infected tooth, removing a foreign body, or relieving bile duct obstruction
    \item \textbf{Pain relief:} paracetamol and anti-inflammatory medicines
    \item \textbf{Wound care:} dressings and packing changes while the cavity heals from within
    \item \textbf{Follow-up:} to confirm resolution and to treat any underlying or predisposing cause
\end{itemize}

\section*{Complications}
\begin{itemize}
    \item Sepsis, a life-threatening systemic response to infection spreading into the bloodstream
    \item Spread of infection to adjacent tissues, such as cellulitis, empyema, or meningitis
    \item Fistula formation, most often after perianal abscesses
    \item Chronic or recurrent abscesses
    \item Scarring and damage to surrounding organs
    \item In brain abscess, raised intracranial pressure and lasting neurological deficit
\end{itemize}

\section*{Prognosis and Outcomes}
Most superficial abscesses resolve within days to about two weeks of drainage, with little long-term effect. Deep abscesses of the liver, lung, or abdomen usually settle with drainage combined with antibiotics, though recovery may take weeks and occasionally requires repeated procedures. Outcomes are generally good when the abscess is drained completely and the underlying cause is addressed. Recurrence is more likely when the predisposing problem, such as diabetes or an ongoing source of infection, remains untreated.

\section*{Prevention and Self-Care}
\begin{itemize}
    \item Clean and cover cuts and scratches promptly
    \item Practise good hand and skin hygiene
    \item Keep up to date with dental care
    \item Manage diabetes and other conditions that weaken immunity
    \item Avoid sharing needles or injecting equipment
    \item Treat skin infections early and do not squeeze boils
    \item After drainage, keep the wound clean and attend follow-up appointments
\end{itemize}

\section*{Sources and Further Reading}
The following organisations provide reliable information on abscesses, their causes, and treatment.

\begin{itemize}
    \item NHS. \textit{Skin abscess: symptoms and treatment.} https://www.nhs.uk/conditions/skin-abscess/
    \item Mayo Clinic. \textit{Boils and carbuncles: overview.} https://www.mayoclinic.org/diseases-conditions/boils/
    \item MSD Manual. \textit{Abscesses: professional overview by site.} https://www.msdmanuals.com/
    \item MedlinePlus. \textit{Abscess health information.} https://medlineplus.gov/abscess.html
    \item Centers for Disease Control and Prevention. \textit{Antimicrobial resistance and skin infection prevention.} https://www.cdc.gov/
    \item World Health Organization. \textit{Infection prevention and control.} https://www.who.int/
\end{itemize}
